% LaTeX Template for AI Problem Analysis Output
% AMC12B 2024 Problem 16 - Strategic Analysis
% Framework Version: 2.1 (Enhanced for COMC)

\documentclass[12pt]{article}
\usepackage[utf8]{inputenc}
\usepackage{amsmath, amssymb, amsthm}
\usepackage{geometry}
\usepackage{xcolor}
\usepackage[most]{tcolorbox}
\usepackage{enumitem}
\usepackage{fancyhdr}

% Page setup
\geometry{margin=1in}
\setlength{\headheight}{14.5pt}
\pagestyle{fancy}
\fancyhf{}
\rhead{AMC12 Problem Analysis}
\lhead{2024 AMC 12B Problem 16}
\cfoot{\thepage}

% Color definitions
\definecolor{problemconfig}{RGB}{230, 242, 255}    % Light blue
\definecolor{strategic}{RGB}{255, 230, 242}       % Light pink
\definecolor{tactical}{RGB}{242, 255, 230}        % Light green
\definecolor{techniques}{RGB}{255, 242, 230}      % Light yellow
\definecolor{verification}{RGB}{255, 230, 230}    % Light red
\definecolor{solution}{RGB}{230, 255, 255}        % Light cyan
\definecolor{competition}{RGB}{242, 242, 255}     % Light purple
\definecolor{gold}{RGB}{218, 165, 32}

% Box styles
\newtcolorbox{problemconfigbox}{
    colback=problemconfig,
    colframe=blue,
    title=PROBLEM CONFIGURATION ANALYSIS,
    fonttitle=\bfseries,
    arc=3pt,
    boxrule=1pt,
    breakable
}

\newtcolorbox{strategicbox}{
    colback=strategic,
    colframe=purple,
    title=STRATEGIC FRAMEWORK APPLICATION,
    fonttitle=\bfseries,
    arc=3pt,
    boxrule=1pt,
    breakable
}

\newtcolorbox{tacticalbox}{
    colback=tactical,
    colframe=green,
    title=TACTICAL METHODOLOGY SELECTION,
    fonttitle=\bfseries,
    arc=3pt,
    boxrule=1pt,
    breakable
}

\newtcolorbox{techniquesbox}{
    colback=techniques,
    colframe=orange,
    title=CONVERSION TECHNIQUES APPLICATION,
    fonttitle=\bfseries,
    arc=3pt,
    boxrule=1pt,
    breakable
}

\newtcolorbox{verificationbox}{
    colback=verification,
    colframe=red,
    title=VERIFICATION STRATEGY DESIGN,
    fonttitle=\bfseries,
    arc=3pt,
    boxrule=1pt,
    breakable
}

\newtcolorbox{solutionbox}{
    colback=solution,
    colframe=cyan,
    title=SOLUTION ROADMAP,
    fonttitle=\bfseries,
    arc=3pt,
    boxrule=1pt,
    breakable
}

\newtcolorbox{competitionbox}{
    colback=competition,
    colframe=violet,
    title=COMPETITION STRATEGY INTEGRATION,
    fonttitle=\bfseries,
    arc=3pt,
    boxrule=1pt,
    breakable
}

\newtcolorbox{keyinsightbox}{
    colback=yellow!20,
    colframe=gold,
    title=KEY INSIGHT,
    fonttitle=\bfseries,
    arc=3pt,
    boxrule=2pt,
    leftrule=5pt,
    breakable
}

\newtcolorbox{warningbox}{
    colback=red!10,
    colframe=red!50,
    title=WARNING: COMMON PITFALL,
    fonttitle=\bfseries,
    arc=3pt,
    boxrule=2pt,
    leftrule=5pt,
    breakable
}

\newtcolorbox{tipbox}{
    colback=green!10,
    colframe=green!50,
    title=PRO TIP,
    fonttitle=\bfseries,
    arc=3pt,
    boxrule=2pt,
    leftrule=5pt,
    breakable
}

\begin{document}

\title{AMC12B 2024 Problem 16\\Strategic Problem Analysis}
\author{Competition Math Framework v2.1}
\date{\today}
\maketitle

\section*{Problem Statement}

A group of $16$ people will be partitioned into $4$ indistinguishable $4$-person committees. Each committee will have one chairperson and one secretary. The number of different ways to make these assignments can be written as $3^{r}M$, where $r$ and $M$ are positive integers and $M$ is not divisible by $3$. What is $r$?

$$\textbf{(A) }5 \qquad \textbf{(B) }6 \qquad \textbf{(C) }7 \qquad \textbf{(D) }8 \qquad \textbf{(E) }9$$

\vspace{0.3cm}
\noindent\textit{Note: The answer is not revealed here to allow students to work through the problem. See the Final Answer section at the end.}

\begin{problemconfigbox}
\subsection*{A. Problem Classification}
\begin{itemize}[leftmargin=*]
    \item \textbf{Problem Type}: To Find (specific numerical value - power of 3)
    \item \textbf{Competition Context}: AMC12B 2024, Problem 16
    \item \textbf{Difficulty Band}: Mid-to-Late band (P16) - Requires sophisticated counting and prime factorization
    \item \textbf{Mathematical Domain}: Combinatorics (counting with restrictions) and Number Theory (prime factorization)
    \item \textbf{Estimated Time}: 4-5 minutes (upper mid-band, requires careful setup)
    \item \textbf{Point Value}: 6 points (standard AMC12 scoring)
\end{itemize}

\subsection*{B. Problem Structure Identification}
\begin{itemize}[leftmargin=*]
    \item \textbf{Given Information}: 
    \begin{itemize}
        \item 16 people to be partitioned
        \item 4 committees, each with exactly 4 people
        \item Committees are \textbf{indistinguishable} (order doesn't matter)
        \item Each committee has 1 chairperson and 1 secretary (distinguishable roles)
        \item The remaining 2 people in each committee have no special role
        \item Total count can be expressed as $3^r M$ where $\gcd(M, 3) = 1$
    \end{itemize}
    
    \item \textbf{Constraints}: 
    \begin{itemize}
        \item Committees are indistinguishable $\Rightarrow$ must divide by $4!$
        \item Each committee must have exactly 4 people
        \item Chairperson and secretary are distinct roles within each committee
        \item Answer must be the exponent $r$ in the factorization
    \end{itemize}
    
    \item \textbf{Objective}: Find the power of 3 (exponent $r$) in the total count
    
    \item \textbf{Hidden Assumptions}: 
    \begin{itemize}
        \item All 16 people are distinguishable
        \item A person cannot serve on multiple committees
        \item Within each committee, chairperson $\neq$ secretary
        \item The two ordinary members of a committee are indistinguishable (no special roles)
    \end{itemize}
    
    \item \textbf{Answer Choice Leverage}: 
    \begin{itemize}
        \item Answers range from 5 to 9
        \item This is a narrow range, suggesting the calculation is reasonably constrained
        \item We're looking for total factors of 3 in numerator minus factors in denominator
        \item Can estimate: $16!$ has about 6 factors of 3, so answer should be close to that
    \end{itemize}
    
    \item \textbf{Proof Requirements}: N/A (computational problem)
\end{itemize}

\subsection*{C. Strategic Orientation}
\begin{itemize}[leftmargin=*]
    \item \textbf{Hypothesis}: The problem requires counting the ways to partition and assign roles, then using Legendre's formula to find factors of 3
    
    \item \textbf{Conclusion}: Need to find $r$ where the total count equals $3^r M$ with $\gcd(M,3)=1$
    
    \item \textbf{Notation}: 
    \begin{itemize}
        \item $N$ = total number of ways to make assignments
        \item $r$ = power of 3 in the prime factorization of $N$
        \item Use Legendre's formula: $\nu_p(n!) = \sum_{k=1}^{\infty} \lfloor n/p^k \rfloor$
    \end{itemize}
    
    \item \textbf{Intuitive Approach}: 
    \begin{enumerate}
        \item Count ways to partition 16 people into 4 groups of 4
        \item Account for indistinguishability of committees (divide by $4!$)
        \item Count ways to assign chairperson and secretary in each committee
        \item Apply Legendre's formula to find total factors of 3
    \end{enumerate}
    
    \item \textbf{Cross-Domain Potential}: 
    \begin{itemize}
        \item This is primarily combinatorics
        \item Requires number theory (prime factorization via Legendre's formula)
        \item Could use multinomial coefficients or sequential selection
    \end{itemize}
\end{itemize}
\end{problemconfigbox}

\begin{strategicbox}
\subsection*{A. Psychological Strategies}
\begin{itemize}[leftmargin=*]
    \item \textbf{Mental Toughness}: This is a multi-step problem that requires both counting skills and number theory. Don't be intimidated - break it into manageable pieces.
    
    \item \textbf{Creativity Cultivation}: Recognize that the "indistinguishable committees" constraint is crucial. Multiple approaches exist (multinomial, sequential selection, arrangement then partition).
    
    \item \textbf{Iterative Refinement}: If the first counting approach seems messy, try an alternative formulation. The key is getting the expression, then applying Legendre's formula.
\end{itemize}

\subsection*{B. Investigation Strategies}
\begin{itemize}[leftmargin=*]
    \item \textbf{Startup Orientation}: Recognize this is a two-part problem:
    \begin{enumerate}
        \item Count the arrangements (combinatorics)
        \item Find factors of 3 (number theory)
    \end{enumerate}
    
    \item \textbf{Get Your Hands Dirty}: 
    \begin{itemize}
        \item Try a smaller case: 4 people, 2 committees of 2 each
        \item Work out the formula structure before worrying about factors of 3
        \item Remember: multinomial $\binom{n}{k_1,k_2,\ldots} = \frac{n!}{k_1!k_2!\cdots}$
    \end{itemize}
    
    \item \textbf{Penultimate Step}: The penultimate step is obtaining the expression for total count. The final step is applying Legendre's formula to count factors of 3.
    
    \item \textbf{Make It Easier}: Focus on getting the counting expression first (even if it's messy). Then systematically apply Legendre's formula to find factors of 3.
    
    \item \textbf{Wishful Thinking}: "What if I could just count factors of 3 in factorials?" You can! That's Legendre's formula.
    
    \item \textbf{Visualization}: Picture dividing 16 people into 4 groups, then assigning 2 special roles in each group.
\end{itemize}

\subsection*{C. Late-Band Strategic Playbook}
\begin{itemize}[leftmargin=*]
    \item \textbf{Answer-Choice Leverage}: The answers cluster around 5-9. Use Legendre on $16!$ to get $\lfloor 16/3 \rfloor + \lfloor 16/9 \rfloor = 5 + 1 = 6$ as a baseline.
    
    \item \textbf{Pattern Recognition}: This is a standard "partition with roles" problem. The formula will have the form $\frac{16!}{\text{symmetry factors}} \times (\text{role assignments})$.
    
    \item \textbf{Verification via Bounds}: The answer should be close to the number of factors of 3 in $16!$, which is 6.
\end{itemize}

\subsection*{D. Argument Construction Strategy}
\begin{itemize}[leftmargin=*]
    \item \textbf{Direct Calculation}: Set up the counting expression, then use Legendre's formula
    
    \item \textbf{Multi-Method Verification}: Can verify by computing the same count using different approaches (multinomial vs. sequential selection)
\end{itemize}

\subsection*{E. Cross-Domain Translation}
\begin{itemize}[leftmargin=*]
    \item \textbf{Combinatorics to Number Theory}: Once we have the counting expression (combinatorics), we translate to finding prime factors (number theory)
    
    \item \textbf{Formula Recognition}: Recognize when to use multinomial coefficients vs. sequential selection
\end{itemize}
\end{strategicbox}

\begin{tacticalbox}
\subsection*{A. Core Mathematical Tactics}
\begin{itemize}[leftmargin=*]
    \item \textbf{Primary Tactics}: 
    \begin{itemize}
        \item \textbf{Symmetry Exploitation}: Account for indistinguishable committees by dividing by $4!$
        \item \textbf{Sequential Counting}: Build up the count step-by-step
        \item \textbf{Prime Factorization}: Use Legendre's formula to count factors of a prime
    \end{itemize}
    
    \item \textbf{Domain-Specific Tactics}:
    \begin{itemize}
        \item \textbf{Combinatorics}:
        \begin{itemize}
            \item Multinomial coefficient: $\binom{n}{k_1,k_2,\ldots,k_r} = \frac{n!}{k_1!k_2!\cdots k_r!}$
            \item For indistinguishable groups, divide by the number of group permutations
            \item Ordered vs. unordered selections
        \end{itemize}
        
        \item \textbf{Number Theory}:
        \begin{itemize}
            \item \textbf{Legendre's Formula}: $\nu_p(n!) = \sum_{k=1}^{\infty} \lfloor n/p^k \rfloor$
            \item For $p=3$: $\nu_3(n!) = \lfloor n/3 \rfloor + \lfloor n/9 \rfloor + \lfloor n/27 \rfloor + \cdots$
            \item Factor counting: $\nu_p(a \cdot b) = \nu_p(a) + \nu_p(b)$
            \item Factor counting: $\nu_p(a / b) = \nu_p(a) - \nu_p(b)$
        \end{itemize}
    \end{itemize}
    
    \item \textbf{Crossover Tactics}: This problem requires both combinatorial counting and number-theoretic factorization
\end{itemize}
\end{tacticalbox}

\begin{techniquesbox}
\subsection*{A. High-Probability Techniques}
\begin{itemize}[leftmargin=*]
    \item \textbf{Primary Technique}: \textbf{Multinomial Coefficients + Legendre's Formula}
    
    \item \textbf{Recognition Cues}: 
    \begin{itemize}
        \item Partition into multiple groups $\Rightarrow$ Multinomial coefficients
        \item "Indistinguishable" groups $\Rightarrow$ Divide by symmetry factor
        \item "Power of prime $p$" $\Rightarrow$ Legendre's formula
        \item Answer asks for exponent $\Rightarrow$ Count prime factors
    \end{itemize}
    
    \item \textbf{Application}:
    \begin{enumerate}
        \item Express total count using multinomial or sequential selection
        \item Simplify to form $\frac{A!}{B! \cdot C! \cdots}$
        \item Apply Legendre's formula to count factors of 3 in numerator and denominator
        \item Compute: $r = \nu_3(\text{numerator}) - \nu_3(\text{denominator})$
    \end{enumerate}
\end{itemize}

\subsection*{B. Alternative Techniques}
\begin{itemize}[leftmargin=*]
    \item \textbf{Alternative Approach 1}: \textbf{Sequential Selection}
    \begin{itemize}
        \item Pick chairperson, secretary, and 2 others for each committee sequentially
        \item More steps but conceptually straightforward
        \item Formula: $\frac{16 \cdot 15 \cdot \binom{14}{2} \cdot 12 \cdot 11 \cdot \binom{10}{2} \cdot 8 \cdot 7 \cdot \binom{6}{2} \cdot 4 \cdot 3 \cdot \binom{2}{2}}{4!}$
    \end{itemize}
    
    \item \textbf{Alternative Approach 2}: \textbf{Arrangement Then Partition}
    \begin{itemize}
        \item Arrange all 16 people in a line: $16!$ ways
        \item Divide into 4 groups of 4: divide by $4!$ for group order
        \item First 2 in each group are chairperson and secretary: divide by $2!$ for each of the 4 groups (since the last 2 are unordered)
        \item Formula: $\frac{16!}{4! \cdot 2^4}$
    \end{itemize}
    
    \item \textbf{Technique Integration}: All approaches should yield the same prime factorization
\end{itemize}
\end{techniquesbox}

\begin{verificationbox}
\subsection*{A. Verification Level Selection}
\begin{itemize}[leftmargin=*]
    \item \textbf{Chosen Level}: Level 5 (Thorough Verification) - 1-2 minutes
    
    \item \textbf{Justification}: 
    \begin{itemize}
        \item Problem 16 is upper mid-band, at the transition to late-band
        \item Multiple counting steps increase error risk
        \item Legendre's formula must be applied correctly
        \item Should verify both the counting expression and the factor calculation
    \end{itemize}
    
    \item \textbf{Time Budget}: 1.5-2 minutes for verification
\end{itemize}

\subsection*{B. Verification Checklist}
\begin{itemize}[leftmargin=*]
    \item \textbf{Counting Expression Verification}:
    \begin{itemize}
        \item Check: Did I account for indistinguishable committees (divide by $4!$)?
        \item Check: Did I correctly count role assignments (chairperson and secretary)?
        \item Check: Did I avoid overcounting or undercounting?
        \item Verify: Does the expression simplify to a clean form?
    \end{itemize}
    
    \item \textbf{Legendre's Formula Application}:
    \begin{itemize}
        \item Check: $\nu_3(16!) = \lfloor 16/3 \rfloor + \lfloor 16/9 \rfloor = 5 + 1 = 6$
        \item Check: $\nu_3(4!) = \lfloor 4/3 \rfloor = 1$
        \item Check: Did I account for all factorial terms in numerator and denominator?
        \item Check: $\nu_3(12^4) = 4 \cdot \nu_3(12) = 4 \cdot 1 = 4$
    \end{itemize}
    
    \item \textbf{Arithmetic Verification}:
    \begin{itemize}
        \item Verify: $r = \nu_3(\text{numerator}) - \nu_3(\text{denominator})$
        \item Double-check addition and subtraction of exponents
        \item Ensure no factors of 3 were missed
    \end{itemize}
    
    \item \textbf{Answer Reasonableness}:
    \begin{itemize}
        \item The answer should be close to $\nu_3(16!) = 6$
        \item Check if answer matches one of the five choices
        \item Verify the arithmetic one more time if answer seems off
    \end{itemize}
\end{itemize}

\subsection*{C. Domain-Specific Verification}
\begin{itemize}[leftmargin=*]
    \item \textbf{Combinatorics Verification}:
    \begin{itemize}
        \item Small case check: For 4 people, 2 committees of 2, does the formula make sense?
        \item Symmetry check: Did I correctly handle indistinguishable objects?
    \end{itemize}
    
    \item \textbf{Number Theory Verification}:
    \begin{itemize}
        \item Legendre verification: Recalculate $\lfloor n/3^k \rfloor$ for relevant $n$ values
        \item Alternative method: Try different counting approach and verify same prime factorization
    \end{itemize}
\end{itemize}
\end{verificationbox}

\begin{keyinsightbox}
\textbf{The Critical Insights}:

\begin{enumerate}
    \item \textbf{Indistinguishability matters}: Since committees are indistinguishable, we must divide by $4!$ to avoid overcounting arrangements of the same 4 committees.
    
    \item \textbf{Two-stage process}: 
    \begin{itemize}
        \item Stage 1: Partition 16 people into 4 groups of 4
        \item Stage 2: Assign chairperson and secretary roles within each group
    \end{itemize}
    
    \item \textbf{Legendre's Formula is essential}: To find the power of 3, we need $\nu_3(n!) = \lfloor n/3 \rfloor + \lfloor n/9 \rfloor + \lfloor n/27 \rfloor + \cdots$
    
    \item \textbf{The clean formula}: The total count equals $\frac{16!}{(4!)^5} \cdot 12^4$, which simplifies the factorization:
    \begin{itemize}
        \item Numerator: $16!$ has 6 factors of 3, $12^4$ has 4 factors of 3
        \item Denominator: $(4!)^5$ has 5 factors of 3
        \item Result: $r = 6 + 4 - 5 = 5$
    \end{itemize}
\end{enumerate}
\end{keyinsightbox}

\begin{solutionbox}
\subsection*{A. Step-by-Step Solution}

\textbf{Step 1: Set Up the Counting Expression}

We need to count two things:
\begin{enumerate}
    \item Ways to partition 16 people into 4 indistinguishable groups of 4
    \item Ways to assign chairperson and secretary roles within each group
\end{enumerate}

\textbf{Checkpoint}: Understanding the structure is crucial before calculating.

\vspace{0.3cm}

\textbf{Step 2: Count Partitions into Committees}

To partition 16 people into 4 distinguishable groups of 4:
\begin{itemize}
    \item Choose 4 people for first committee: $\binom{16}{4}$
    \item Choose 4 people for second committee: $\binom{12}{4}$
    \item Choose 4 people for third committee: $\binom{8}{4}$
    \item Remaining 4 people form fourth committee: $\binom{4}{4} = 1$
\end{itemize}

Total for distinguishable committees: $\binom{16}{4} \cdot \binom{12}{4} \cdot \binom{8}{4} \cdot \binom{4}{4}$

Since committees are \textbf{indistinguishable}, we divide by $4!$:

$$\text{Partitions} = \frac{1}{4!} \cdot \binom{16}{4} \cdot \binom{12}{4} \cdot \binom{8}{4} \cdot \binom{4}{4}$$

This simplifies to:
\begin{align*}
    \text{Partitions} &= \frac{1}{4!} \cdot \frac{16!}{4! \cdot 12!} \cdot \frac{12!}{4! \cdot 8!} \cdot \frac{8!}{4! \cdot 4!} \cdot \frac{4!}{4! \cdot 0!} \\
    &= \frac{16!}{4! \cdot (4!)^4} \\
    &= \frac{16!}{(4!)^5}
\end{align*}

\textbf{Checkpoint}: Does this formula make sense? Yes - we have $16!$ for ordering, divided by $(4!)^4$ for internal committee orderings, divided by $4!$ for committee order.

\vspace{0.3cm}

\textbf{Step 3: Count Role Assignments}

Within each committee of 4 people:
\begin{itemize}
    \item Choose chairperson: 4 ways
    \item Choose secretary (from remaining 3): 3 ways
    \item Total per committee: $4 \times 3 = 12$ ways
\end{itemize}

For all 4 committees: $12^4$ ways

\vspace{0.3cm}

\textbf{Step 4: Compute Total Count}

Total number of ways:
$$N = \frac{16!}{(4!)^5} \cdot 12^4$$

\textbf{Checkpoint}: This is our counting expression. Now we need to find the power of 3 in $N$.

\vspace{0.3cm}

\textbf{Step 5: Apply Legendre's Formula to Find Factors of 3}

\textbf{Legendre's Formula}: $\nu_3(n!) = \sum_{k=1}^{\infty} \lfloor n/3^k \rfloor$

\textbf{For $16!$}:
\begin{align*}
    \nu_3(16!) &= \lfloor 16/3 \rfloor + \lfloor 16/9 \rfloor + \lfloor 16/27 \rfloor + \cdots \\
    &= 5 + 1 + 0 + \cdots \\
    &= 6
\end{align*}

\textbf{For $(4!)^5$}:
\begin{align*}
    \nu_3(4!) &= \lfloor 4/3 \rfloor + \lfloor 4/9 \rfloor + \cdots \\
    &= 1 + 0 + \cdots \\
    &= 1
\end{align*}

Therefore: $\nu_3((4!)^5) = 5 \cdot \nu_3(4!) = 5 \cdot 1 = 5$

\textbf{For $12^4$}:
\begin{align*}
    12 &= 4 \times 3 = 2^2 \times 3 \\
    \nu_3(12) &= 1
\end{align*}

Therefore: $\nu_3(12^4) = 4 \cdot \nu_3(12) = 4 \cdot 1 = 4$

\vspace{0.3cm}

\textbf{Step 6: Calculate the Total Power of 3}

The total count is:
$$N = \frac{16! \cdot 12^4}{(4!)^5}$$

Therefore:
\begin{align*}
    r = \nu_3(N) &= \nu_3(16!) + \nu_3(12^4) - \nu_3((4!)^5) \\
    &= 6 + 4 - 5 \\
    &= 5
\end{align*}

\textbf{Final Answer}: $r = 5$, which is choice \textbf{(A)}.

\subsection*{B. Verification of Solution}

\textbf{Verification 1: Check Legendre Calculations}

\begin{itemize}
    \item $\nu_3(16!)$: $\lfloor 16/3 \rfloor = 5$, $\lfloor 16/9 \rfloor = 1$, $\lfloor 16/27 \rfloor = 0$ $\Rightarrow$ $5+1 = 6$ \checkmark
    \item $\nu_3(4!)$: $\lfloor 4/3 \rfloor = 1$, $\lfloor 4/9 \rfloor = 0$ $\Rightarrow$ $1$ \checkmark
    \item $\nu_3((4!)^5) = 5 \times 1 = 5$ \checkmark
    \item $\nu_3(12) = 1$ since $12 = 2^2 \times 3$ \checkmark
    \item $\nu_3(12^4) = 4 \times 1 = 4$ \checkmark
\end{itemize}

\textbf{Verification 2: Final Arithmetic}
$$r = 6 + 4 - 5 = 5$$ \checkmark

\textbf{Verification 3: Reasonableness Check}

The answer $r=5$ is close to $\nu_3(16!) = 6$, which makes sense because:
\begin{itemize}
    \item We're dividing by $(4!)^5$ which removes 5 factors of 3
    \item We're multiplying by $12^4$ which adds 4 factors of 3
    \item Net effect: $6 - 5 + 4 = 5$ \checkmark
\end{itemize}

\subsection*{C. Alternative Approach: Sequential Selection}

We can verify our answer using a different counting method.

Pick members sequentially for each committee:
\begin{itemize}
    \item Committee 1: Pick chairperson (16 ways), secretary (15 ways), two others $\binom{14}{2}$ ways
    \item Committee 2: Pick chairperson (12 ways), secretary (11 ways), two others $\binom{10}{2}$ ways
    \item Committee 3: Pick chairperson (8 ways), secretary (7 ways), two others $\binom{6}{2}$ ways
    \item Committee 4: Pick chairperson (4 ways), secretary (3 ways), two others $\binom{2}{2}$ ways
\end{itemize}

Divide by $4!$ since committees are indistinguishable:

$$N = \frac{16 \cdot 15 \cdot \binom{14}{2} \cdot 12 \cdot 11 \cdot \binom{10}{2} \cdot 8 \cdot 7 \cdot \binom{6}{2} \cdot 4 \cdot 3 \cdot \binom{2}{2}}{4!}$$

Expanding binomial coefficients:
\begin{align*}
    N &= \frac{16 \cdot 15 \cdot \frac{14 \cdot 13}{2} \cdot 12 \cdot 11 \cdot \frac{10 \cdot 9}{2} \cdot 8 \cdot 7 \cdot \frac{6 \cdot 5}{2} \cdot 4 \cdot 3 \cdot 1}{4!} \\
    &= \frac{16! / 2^4}{4!} \\
    &= \frac{16!}{4! \cdot 2^4}
\end{align*}

This is equivalent to our previous expression! To verify:
\begin{align*}
    \frac{16!}{4! \cdot 2^4} &= \frac{16!}{(4!)^5} \cdot 12^4
\end{align*}

Since $12^4 = (4 \cdot 3)^4 = 2^8 \cdot 3^4$ and $(4!)^4 = (2^3 \cdot 3)^4 = 2^{12} \cdot 3^4$, we have:
\begin{align*}
    \frac{16! \cdot 12^4}{(4!)^5} &= \frac{16! \cdot 2^8 \cdot 3^4}{(4!)^5} \\
    &= \frac{16!}{4! \cdot (4!)^4} \cdot 2^8 \cdot 3^4 \\
    &= \frac{16!}{4! \cdot 2^{12} \cdot 3^4} \cdot 2^8 \cdot 3^4 \\
    &= \frac{16!}{4! \cdot 2^4}
\end{align*}

Both give $r = 6 + 4 - 5 = 5$ \checkmark

\end{solutionbox}

\begin{competitionbox}
\subsection*{A. Time Management}
\begin{itemize}[leftmargin=*]
    \item \textbf{Estimated Time}: 4-5 minutes total
    \begin{itemize}
        \item Problem comprehension: 45 seconds
        \item Set up counting expression: 1.5 minutes
        \item Apply Legendre's formula: 1.5 minutes
        \item Calculate final answer: 30 seconds
        \item Verification: 1 minute
    \end{itemize}
    
    \item \textbf{Time Checkpoints}: 
    \begin{itemize}
        \item At 1 minute: Should understand the problem structure
        \item At 2.5 minutes: Should have counting expression
        \item At 4 minutes: Should have applied Legendre's formula
        \item At 5 minutes: Should be done with verification
    \end{itemize}
    
    \item \textbf{Priority Level}: \textbf{High-Medium} - This is P16, at the boundary of mid/late band. Accessible if you know Legendre's formula, but can be time-consuming.
\end{itemize}

\subsection*{B. Risk Assessment}
\begin{itemize}[leftmargin=*]
    \item \textbf{Confidence Level}: 85-90\% after verification
    \begin{itemize}
        \item The counting structure is standard
        \item Legendre's formula is mechanical
        \item Alternative method confirms the result
        \item Answer matches expected range
    \end{itemize}
    
    \item \textbf{Common Error Sources}:
    \begin{itemize}
        \item Forgetting to divide by $4!$ for indistinguishable committees
        \item Incorrectly counting role assignments
        \item Arithmetic errors in Legendre's formula
        \item Forgetting factors of 3 in $12^4$
    \end{itemize}
    
    \item \textbf{Abandonment Criteria}: 
    \begin{itemize}
        \item If you don't know Legendre's formula $\Rightarrow$ skip this problem
        \item If stuck on counting after 3 minutes $\Rightarrow$ make educated guess
        \item If calculation is taking > 6 minutes $\Rightarrow$ flag and move on
    \end{itemize}
    
    \item \textbf{Flag for Review}: Yes if time permits - this is a problem where small arithmetic errors can occur
\end{itemize}

\subsection*{C. Strategic Context}
\begin{itemize}[leftmargin=*]
    \item \textbf{Problem Positioning}: P16 is the transition point to late-band. It's worth attempting if you're comfortable with combinatorics and number theory.
    
    \item \textbf{Score Impact}: Worth 6 points. Getting this right demonstrates strong technical skills.
    
    \item \textbf{Time Budget Implications}: Spending 4-5 minutes here is reasonable if you have a clear path to the solution.
    
    \item \textbf{Technique Practice}: This problem tests:
    \begin{itemize}
        \item Partition counting with indistinguishability
        \item Role assignment counting
        \item Legendre's formula for prime factorization
        \item Multi-step problem organization
    \end{itemize}
\end{itemize}
\end{competitionbox}

\begin{warningbox}
\textbf{Common Pitfalls to Avoid}:

\begin{enumerate}
    \item \textbf{Forgetting indistinguishability}: The committees are indistinguishable, so you MUST divide by $4!$. This is easy to miss.
    
    \item \textbf{Incorrect role counting}: Each committee has $4 \times 3 = 12$ ways to assign chairperson and secretary, not $\binom{4}{2}$. Order matters for these roles!
    
    \item \textbf{Legendre's formula errors}:
    \begin{itemize}
        \item Remember to include all powers: $\lfloor n/3 \rfloor + \lfloor n/9 \rfloor + \lfloor n/27 \rfloor + \cdots$
        \item For $n=16$: Don't forget $\lfloor 16/9 \rfloor = 1$
    \end{itemize}
    
    \item \textbf{Forgetting $12^4$ factors}: $12 = 2^2 \times 3$, so $12^4$ contributes 4 factors of 3 to the numerator.
    
    \item \textbf{Sign errors in exponent calculation}: When dividing, subtract exponents: $\nu_3(A/B) = \nu_3(A) - \nu_3(B)$
    
    \item \textbf{Arithmetic errors}: This problem has multiple calculations. Double-check: $6 + 4 - 5 = 5$.
    
    \item \textbf{Confusing distinguishable vs. indistinguishable}: Make sure you understand what's being counted.
\end{enumerate}
\end{warningbox}

\begin{tipbox}
\textbf{Competition Tips}:

\begin{enumerate}
    \item \textbf{Know Legendre's Formula}: This is essential for late AMC12 problems. Memorize it:
    $$\nu_p(n!) = \sum_{k=1}^{\infty} \lfloor n/p^k \rfloor$$
    
    \item \textbf{Quick Legendre for small $n$}: For $n \leq 27$ and $p=3$, you only need:
    $$\nu_3(n!) = \lfloor n/3 \rfloor + \lfloor n/9 \rfloor + \lfloor n/27 \rfloor$$
    
    \item \textbf{Recognize standard partition problems}: "Partition into indistinguishable groups" is a classic setup. Formula: $\frac{n!}{(k!)^m \cdot m!}$ for $n$ objects into $m$ indistinguishable groups of $k$ each.
    
    \item \textbf{Role counting shortcut}: For $r$ distinguishable roles chosen from $n$ people: $n(n-1)(n-2)\cdots(n-r+1) = \frac{n!}{(n-r)!}$
    
    \item \textbf{Answer checking}: The answer should be close to $\nu_3(16!) = 6$. If you get 15 or -3, you made an error!
    
    \item \textbf{Use answer choices}: If stuck, eliminate impossible answers:
    \begin{itemize}
        \item Too high (> 7)? Probably overcounted factors
        \item Too low (< 4)? Probably undercounted factors
    \end{itemize}
    
    \item \textbf{Simplify before factorizing}: Get the cleanest expression first, then apply Legendre's formula. This reduces arithmetic errors.
\end{enumerate}
\end{tipbox}

\section*{Final Answer}

The total number of ways to make the assignments is:
$$N = \frac{16!}{(4!)^5} \cdot 12^4$$

Using Legendre's formula to count factors of 3:
\begin{align*}
    \nu_3(16!) &= 6 \\
    \nu_3((4!)^5) &= 5 \\
    \nu_3(12^4) &= 4
\end{align*}

Therefore:
$$r = \nu_3(N) = 6 + 4 - 5 = 5$$

$$\boxed{\textbf{(A) } 5}$$

\section*{Confidence Assessment}
\begin{itemize}[leftmargin=*]
    \item \textbf{Confidence Level}: 90\%
    \begin{itemize}
        \item Counting expression is standard and verified
        \item Legendre's formula calculations are straightforward
        \item Alternative sequential method confirms the same result
        \item Answer is in the expected range (close to 6)
        \item All arithmetic double-checked
    \end{itemize}
    
    \item \textbf{Verification Applied}: Level 5 (Thorough Verification) plus alternative method
    \begin{itemize}
        \item Verified counting expression
        \item Rechecked all Legendre calculations
        \item Confirmed arithmetic: $6 + 4 - 5 = 5$
        \item Cross-verified using sequential selection approach
        \item Verified answer matches a given choice
    \end{itemize}
    
    \item \textbf{Flag for Review}: Maybe - if extra time is available, worth a quick recheck of arithmetic
    
    \item \textbf{Time Spent}: Approximately 4.5 minutes (within target range)
    \begin{itemize}
        \item 45 seconds: Problem comprehension
        \item 1.5 minutes: Setting up counting expression
        \item 1.5 minutes: Applying Legendre's formula
        \item 30 seconds: Final calculation
        \item 30 seconds: Quick verification
    \end{itemize}
\end{itemize}

\section*{Learning Points}

\textbf{Key Takeaways from This Problem}:

\begin{enumerate}
    \item \textbf{Legendre's Formula Mastery}: This is an essential tool for counting prime factors in factorials. Know it cold.
    
    \item \textbf{Indistinguishability Matters}: When objects (committees) are indistinguishable, divide by the symmetry factor (here, $4!$).
    
    \item \textbf{Multi-Stage Counting}: Break complex counting problems into stages:
    \begin{itemize}
        \item Stage 1: Partition into groups
        \item Stage 2: Assign roles within groups
    \end{itemize}
    
    \item \textbf{Alternative Methods for Verification}: Using different counting approaches (multinomial vs. sequential) provides strong confidence.
    
    \item \textbf{Factor Arithmetic}: When computing $\nu_p(A \cdot B / C)$:
    $$\nu_p(A \cdot B / C) = \nu_p(A) + \nu_p(B) - \nu_p(C)$$
    
    \item \textbf{Simplify First, Factorize Second}: Get the expression in simplest form before worrying about prime factorization.
\end{enumerate}

\vspace{0.5cm}

\textbf{Related Problems to Practice}:
\begin{itemize}
    \item Other AMC12 problems involving Legendre's formula
    \item Partition problems with indistinguishable groups
    \item Combinatorial problems requiring role assignments
    \item Number theory problems involving prime factorization
\end{itemize}

\vspace{0.5cm}

\textbf{Formula to Remember}:

\textbf{Legendre's Formula for prime $p$}:
$$\nu_p(n!) = \sum_{k=1}^{\infty} \left\lfloor \frac{n}{p^k} \right\rfloor$$

For $p=3$ and $n \leq 27$:
$$\nu_3(n!) = \left\lfloor \frac{n}{3} \right\rfloor + \left\lfloor \frac{n}{9} \right\rfloor + \left\lfloor \frac{n}{27} \right\rfloor$$

\end{document}

